% !TEX program = xelatex
\documentclass[12pt,a4paper]{article}

% --- اللغة العربية والخطوط (بسيطة، بدون bidi=basic، مع bidi افتراضي) ---
\usepackage[english, bidi=default]{babel}
\babelprovide[import, main]{arabic}
\babelprovide[import, onchar=ids fonts]{english}
\usepackage{fontspec}
\defaultfontfeatures{Ligatures=TeX}

% خطوط عربية متوفرة في معظم أنظمة ويندوز
\babelfont[arabic]{rm}[Scale=1.1, Script=Arabic, Language=Arabic]{Times New Roman}
\babelfont[arabic]{sf}[Scale=1.1, Script=Arabic, Language=Arabic]{Arial}
\babelfont[arabic]{tt}[Scale=1.1, Script=Arabic, Language=Arabic]{Courier New}
% خطوط لاتينية للرياضيات والنصوص الإنجليزية
\babelfont[english]{rm}{Times New Roman}
\babelfont[english]{sf}{Arial}
\babelfont[english]{tt}{Courier New}

% --- حزم الرياضيات الأساسية (بدون booktabs لتجنب التعارض مع bidi) ---
\usepackage{amsmath,amssymb,amsfonts,mathtools}
\usepackage{geometry}
\geometry{left=1.0cm,right=1.0cm,top=1.0cm,bottom=3.4cm}
\usepackage{hyperref}
\usepackage{url}
\usepackage{xcolor}
\usepackage{titlesec}
\usepackage{fancyhdr}
\usepackage{array}

% --- ألوان YUCT ---
\definecolor{skyblue}{RGB}{0,102,204}
\definecolor{darkblue}{RGB}{0,0,139}
\definecolor{gold}{RGB}{255,215,0}

% --- تنسيق الأقسام ---
\titleformat{\section}{\LARGE\bfseries\color{darkblue}}{\thesection}{1em}{}
\titleformat{\subsection}{\Large\bfseries\color{darkblue}}{\thesubsection}{1em}{}

% --- رأس وتذييل الصفحة ---
\fancypagestyle{firstpage}{%
	\fancyhf{}%
	\renewcommand{\headrulewidth}{-5pt}%
	\renewcommand{\footrulewidth}{-5pt}%
	\fancyfoot[C]{%
		\begin{minipage}[t]{\textwidth}
			\centering
			\textcolor{gold}{\rule{\textwidth}{1.6pt}}\\[5pt]
			{\small \textsuperscript{1}Yakushev Research, \url{https://yuct.org/}} \hfill
			{\small بروتوكول YPSDC: \url{https://ypsdc.com/}}\\[-2pt]
			\textcolor{gold}{\rule{\textwidth}{1.6pt}}\\[5pt]
			\textbf{نظرية التنسيق الموحدة ياكوشيف 2026} \hfill \thepage\\[10pt]
		\end{minipage}%
	}%
}

\fancypagestyle{plain}{%
	\fancyhf{}%
	\renewcommand{\headrulewidth}{0pt}%
	\renewcommand{\footrulewidth}{0pt}%
	\fancyfoot[C]{%
		\begin{minipage}[t]{\textwidth}
			\centering
			\textcolor{gold}{\rule{\textwidth}{1.6pt}}\\[4pt]
			\textbf{نظرية التنسيق الموحدة ياكوشيف 2026} \hfill \thepage\\[4pt]
		\end{minipage}%
	}%
}

\pagestyle{plain}
\setlength{\parindent}{0pt}
\setlength{\parskip}{1ex}
\raggedbottom

\hypersetup{
	colorlinks=true,
	linkcolor=blue,
	citecolor=blue,
	urlcolor=blue,
}

% --- رأس الصفحة (شعار YUCT) ---
\newcommand{\makeheader}{%
	\noindent
	\begin{minipage}{0.17\textwidth}
		\raggedright
		{\sffamily\bfseries\color{skyblue}\fontsize{46}{46}\selectfont YUCT}%
	\end{minipage}%
	\hspace{30pt}
	\begin{minipage}{0.32\textwidth}
		\raggedright
		{\sffamily\fontsize{11}{13}\selectfont نظرية\\ التنسيق\\ الموحدة ياكوشيف}%
	\end{minipage}%
	\hfill
	\begin{minipage}{0.38\textwidth}
		\raggedleft
		{\sffamily\bfseries فرضية / خارطة طريق}\\ 
		{\sffamily الإصدار 2.0 / يونيو 2026}\\ 
	\end{minipage}
	\par\vspace{5pt}
	\noindent\textcolor{gold}{\rule{\textwidth}{1.6pt}}
	\par\vspace{0pt}
}

\begin{document}
	\thispagestyle{firstpage}
	\makeheader
	
	\vspace{0cm}
	{\LARGE\bfseries مشكلة CP القوية:\\[0.1cm] فرضية جبرية محددة ضمن YUCT\\[0.2cm] \Large خارطة طريق للتحقق النظري والتجريبي \par}
	\vspace{0.2cm}
	
	{\large أليكسي ف. ياكوشيف\textsuperscript{1}} \\
	\vspace{0.0cm}
	
	{\color{darkblue}\textbf{\LARGE الملخص}} \par
	{\large
		\noindent
		مشكلة CP القوية – لماذا معامل \(\theta\) المخالف لـ CP في ديناميكا الكم اللونية (QCD) أصغر من \(10^{-10}\) – هي اللغز الأخير غير المحلول للنموذج العياري ضمن نظرية التنسيق الموحدة ياكوشيف (YUCT). في هذه الوثيقة نقدم فرضية جبرية ملموسة تنتج طبيعياً الكبت المطلوب. نفترض أن عمق التنسيق الفعال لمعامل \(\theta\) هو
		\[
		L_\theta = \frac{1}{2}\,L_W \;+\; \frac{S_{\mathrm{odd}}}{S_{\mathrm{even}}},
		\]
		حيث \(L_W = 96\) هو عمق بوزون \(W\) و \(S_{\mathrm{odd}}/S_{\mathrm{even}} = 3/2\) هو ثابت أساسي للنظرية. هذه الصيغة لا تحتوي على بارامترات حرة وتعطي \(L_\theta = 49.5\)، وهو ما يقابل \(\theta \sim (2/3)^{49.5} \approx 2.3 \times 10^{-10}\). نناقش الدافع الفيزيائي لهذا التعبير، ونضمّنه في الإطار الجبري الأوسع لـ YUCT، ونحدد برنامجاً من الاختبارات التجريبية – بشكل أساسي من خلال عزم ثنائي القطب الكهربائي للنيوترون (nEDM) – يمكن أن يدحض الفرضية. يُقدَّم الاقتراح كتخمين واضح ودقيق رياضياً، مفتوح للدحض أو التأكيد بواسطة البيانات المستقبلية.
	}
	
	\vspace{0.1cm}
	\noindent
	{\color{darkblue}\textbf{الكلمات المفتاحية:}} YUCT, مشكلة CP القوية, فرضية جبرية, عمق التنسيق, عزم ثنائي القطب الكهربائي للنيوترون, برنامج بحثي مفتوح.
	\vspace{0.4cm}
	
	\clearpage
	\tableofcontents
	\newpage
	
	\section{مقدمة}
	\label{sec:intro}
	
	قدمت نظرية التنسيق الموحدة ياكوشيف (YUCT) حلولاً كمية خالية من البارامترات لمشكلة التسلسل الهرمي، ومشكلة الثابت الكوني، ونمط كتل الفرميونات، والعديد من الألغاز الكبرى الأخرى في الفيزياء الأساسية~\cite{Y1,Y3,Y4}. الاستثناء الوحيد كان مشكلة CP القوية: الصغر غير الطبيعي لمعامل \(\theta\) المخالف لـ CP في ديناميكا الكم اللونية (QCD)، والذي تم تقييده تجريبياً بـ \(|\theta| \lesssim 10^{-10}\)~\cite{nEDM2020}.
	
	المحاولات السابقة داخل YUCT لتعيين عمق تنسيق لـ \(\theta\) بواسطة عدد صحيح وهمي \(k_\theta\) كانت غير مرضية، لأنها اختزلت إلى ملاءمة بدلاً من تنبؤ بنيوي. في هذه المذكرة نقدم نهجاً مختلفاً: نستغل التفاعل بين القطاع المنتظم (الفرميوني) والقطاع الفوضوي (الطوبولوجي) لفراغ QCD – وهو تفاعل تتمتع YUCT بموقع فريد لقياسه بفضل وصفها الجبري الموحد للنظام والفوضى~\cite{Y2,Feig}.
	
	لا ندعي اشتقاقاً صارماً من لاغرانجي 19‑البعدي؛ مثل هذا الاشتقاق هو هدف بعيد المدى. بدلاً من ذلك، نذكر تخميناً جبرياً حاداً يتبع بشكل طبيعي من الأرقام الموجودة بالفعل في YUCT، ولا يحتوي على بارامترات قابلة للتعديل، ويقدم تنبؤاً قابلاً للدحض لعزم ثنائي القطب الكهربائي للنيوترون (nEDM) ولكتلة الأكسيون. يُقدَّم التخمين كهدف محدد جيداً للعمل النظري والتجريبي المستقبلي.
	
	\section{ثوابت النظام والفوضى في YUCT}
	\label{sec:oc}
	
	\subsection{النظام}
	قانون الخطأ العالمي
	\[
	\varepsilon = \kappa_c\,\alpha\,K_{\mathrm{eff}}^{-\beta},\qquad \beta = \frac{2}{3},\quad \kappa_c = \frac{1}{3},
	\]
	مع البنية الهرمية للتنسيق، يؤدي إلى حلقة جبرية مغلقة
	\[
	S_{\mathrm{odd}} = \frac{6}{5} = 1.2,\quad S_{\mathrm{even}} = \frac{4}{5} = 0.8,\quad 
	\frac{S_{\mathrm{odd}}}{S_{\mathrm{even}}} = \frac{3}{2} = \frac{1}{\beta}.
	\]
	جميع الكتل الفيزيائية تُعبر عنها بدلالة كتلة بلانك وعمق تنسيق \(L\):
	\[
	m \approx M_{\mathrm{Pl}}\left(\frac{2}{3}\right)^L.
	\]
	على وجه الخصوص، كتلة بوزون \(W\) تقابل \(L_W = 96\)، وهو العدد الكلي لدرجات حرية فرميونات فايل في النموذج العياري (بما في ذلك النيوترينوات اليمنى)~\cite{Y3}.
	
	\subsection{الفوضى}
	ثوابت فيجنباوم العالمية
	\[
	\delta \approx 4.6692,\qquad \alpha \approx 2.5029,
	\]
	تصف طريق مضاعفة الدورة نحو الفوضى. في YUCT هي مرتبطة جبرياً بثوابت النظام عبر النسبة الذهبية \(\varphi = (1+\sqrt{5})/2\) وثابت الدائرة \(\pi\):
	\[
	2\alpha^2 \approx \pi\delta,\qquad
	\pi \approx S_{\mathrm{odd}}\varphi^2 \quad \Longrightarrow \quad 
	2\alpha^2 \approx (S_{\mathrm{odd}}\varphi^2)\,\delta .
	\]
	هذا الربط يظهر أن ثوابت فيجنباوم ليست غريبة عن YUCT؛ إنها تنتمي إلى نفس النظام الجبري الموسع~\cite{Y2}.
	
	\subsection{الفراغ الكمي كمزدوج النظام-الفوضى}
	فراغ QCD هو الساحة التي تتعايش فيها الجوانب المنتظمة (تكاثف الكواركات، كسر التناظر الكيرالي) والفوضوية (تقلبات الانستنتون، القابلية الطوبولوجية). معامل \(\theta\) يرتبط بكثافة الشحنة الطوبولوجية \(G\widetilde{G}\)، التي تتلقى مساهمات من كلا القطاعين. لذلك من الطبيعي أن نتوقع أن العمق الفعال \(L_\theta\) هو دالة للأعماق التي تميز المكونات المنتظمة والفوضوية.
	
	في الأعمال السابقة قدمنا أعماقاً أولية \(L_{\mathrm{ord}}\) (القطاع الفرميوني) و \(L_{\mathrm{ch}}\) (القطاع الطوبولوجي/الفوضوي)، لكن التركيبات الخطية البسيطة لم تعيد إنتاج الكبت المطلوب \(|\theta| \sim 10^{-10}\). هذا دفعنا للبحث عن تركيبة أكثر بنيوية.
	
	\section{الفرضية الجبرية لـ \(L_\theta\)}
	\label{sec:hypothesis}
	
	\subsection{بيان الفرضية}
	
	نقترح أن عمق التنسيق الفعال لمعامل \(\theta\) يُعطى بالصيغة البسيطة
	\begin{equation}
		\boxed{L_\theta = \frac{1}{2}\,L_W \;+\; \frac{S_{\mathrm{odd}}}{S_{\mathrm{even}}}},
		\label{eq:Ltheta}
	\end{equation}
	حيث
	\begin{itemize}
		\item \(L_W = 96\) هو عمق بوزون \(W\) (أساس المقياس الكهروضعيف)؛
		\item \(S_{\mathrm{odd}}/S_{\mathrm{even}} = 3/2 = 1.5\) هو النسبة العالمية التي تحكم عامل الكتلة بين الأجيال وتظهر في الحلقة الجبرية الأولية.
	\end{itemize}
	
	بهذه الأرقام،
	\[
	L_\theta = \frac{96}{2} + 1.5 = 48 + 1.5 = 49.5.
	\]
	
	معامل \(\theta\) المطابق هو، حتى عامل مقداره واحد متوقع من حساب الانستنتون،
	\[
	\theta \sim \left(\frac{2}{3}\right)^{L_\theta} = \left(\frac{2}{3}\right)^{49.5} \approx 2.3 \times 10^{-10}.
	\]
	هذا يقع تماماً داخل النافذة التجريبية المسموحة \(|\theta| \lesssim 10^{-10}\)~\cite{nEDM2020}.
	
	\subsection{التفسير الفيزيائي}
	
	لماذا يجب أن تظهر هذه التركيبة بالذات؟
	\begin{enumerate}
		\item \textbf{العامل \(1/2\)}: المقياس الكهروضعيف يتحدد بقيمة توقع الفراغ \(v \approx 246\) GeV، والتي تُحصل من العمق الأساسي 96 من خلال عامل هندسي \(\xi_W \approx 0.53\). العمق 48 يقابل المقياس \(\sqrt{v M_{\mathrm{Pl}}}\) – الوسط الهندسي بين المقياس الضعيف ومقياس بلانك – وهو بالضبط المقياس الذي تتوقع عنده المساهمات الفرميونية والبوزونية في طاقة الفراغ أن تتقاطع. في YUCT، من المعروف أن هذا المقياس يلعب دوراً خاصاً في ديناميكيات تنسيق مكثف هيغز (ملحق~\(\Lambda\)).
		\item \textbf{الإزاحة \(S_{\mathrm{odd}}/S_{\mathrm{even}} = 3/2\)}: النسبة \(3/2 = 1/\beta\) تحدد اللامتماثل غير القابل للاختزال بين الارتباطات المنتظمة (المستويات الفردية) وغير المنتظمة (المستويات الزوجية) في أي نظام تنسيق هرمي. وجودها في \(L_\theta\) يشير إلى أن الحد المخالف لـ CP ينشأ من اختلال صغير لا مفر منه بين القطاع المهيمن بالنظام والقطاع المهيمن بالفوضى للفراغ. أن هذه الإزاحة مضافة وليست مضاعفة يعكس البنية المنفصلة ذات الطبقات لشبكة التنسيق.
	\end{enumerate}
	
	لذلك يمكن قراءة المعادلة~(\ref{eq:Ltheta}) على النحو التالي: عمق \(\theta\) هو نصف العمق الضعيف مضافاً إليه الإزاحة العالمية للنظام-الفوضى. لا تحتوي على بارامترات حرة وتستخدم فقط أرقاماً تم تحديدها بالفعل بشكل مستقل داخل YUCT.
	
	\subsection{نتائج فورية}
	
	\paragraph{عزم ثنائي القطب الكهربائي للنيوترون.}
	nEDM مرتبط بـ \(\theta\) بالعلاقة \(d_n \approx 3.6 \times 10^{-16}\,\theta\; e\!\cdot\!\mathrm{cm}\) (مع عامل عدم يقين هادروني)~\cite{nEDM2020}. بأخذ \(\theta \approx 2.3 \times 10^{-10}\) نحصل على
	\[
	d_n \approx 8.3 \times 10^{-26}\; e\!\cdot\!\mathrm{cm}.
	\]
	الحدود التجريبية الحالية هي \(d_n < 1.8 \times 10^{-26}\; e\!\cdot\!\mathrm{cm}\) (بنسبة ثقة 90\%). الجيل القادم من التجارب في PSI و SNS و TUCAN تستهدف حساسيات تصل إلى بضع \(10^{-28}\; e\!\cdot\!\mathrm{cm}\). نتيجة صفرية عند هذا المستوى ستدحض الفرضية؛ كشف إيجابي سيكون تأكيداً دراماتيكياً.
	
	\paragraph{فيزياء الأكسيون.}
	إذا تم حل مشكلة CP القوية بواسطة آلية بيكي-كوين، فإن كتلة الأكسيون \(m_a\) ترتبط بـ \(\theta\) وبثابت اضمحلال الأكسيون \(f_a\). المعادلة~(\ref{eq:Ltheta}) تحدد \(f_a\) (وبالتالي \(m_a\)) من خلال نفس منطق أعماق التنسيق. هذا يوفر هدفاً تجريبياً تكميلياً لـ ADMX وبحوث الأكسيون الأخرى.
	
	\section{خارطة طريق للتحقق}
	\label{sec:roadmap}
	
	الفرضية دقيقة رياضياً وقابلة للدحض. نقترح البرنامج التالي:
	
	\subsection{المهام النظرية}
	\begin{enumerate}
		\item \textbf{الاشتقاق من لاغرانجي 19‑البعدي.}
		يجب اختزال الفعل الفعال للنظرية الأم 19‑البعدية إلى حد QCD القياسي \(G\widetilde{G}\). الهدف هو إظهار أن المعاملات في الاختزال تنتج طبيعياً العمق \(L_\theta = 49.5\) عند تحديد حالة الفراغ الصحيحة. هذه مشكلة رياضية صعبة ولكنها محددة جيداً.
		\item \textbf{نظرية المقياس الشبكية.}
		حسابات QCD الشبكية الحالية يمكنها حساب القابلية الطوبولوجية واعتماد طاقة الفراغ على \(\theta\). تحليل مستوحى من YUCT يجب أن يختبر ما إذا كانت البنية الجبرية للمعادلة~(\ref{eq:Ltheta}) تنعكس في تحجيم نتائج الشبكة مع ثابت الارتباط، أو في توزيع عناقيد الشحنة الطوبولوجية.
		\item \textbf{تحسين التنبؤ بـ nEDM.}
		عناصر المصفوفة الهاديرونية التي تربط \(\theta\) بـ \(d_n\) معروفة من نظرية الاضطراب الكيرالية و QCD الشبكية. حساب مخصص لهذه العناصر، مقترناً بقيمة YUCT لـ \(\theta\)، سيعطي تنبؤاً حاداً يمكن مقارنته مباشرة بحدود الجيل القادم من تجارب nEDM.
	\end{enumerate}
	
	\subsection{المهام التجريبية}
	\begin{enumerate}
		\item \textbf{تجارب nEDM (PSI, SNS, TUCAN).}
		nEDM مقاس أعلى من \(10^{-26}\;e\!\cdot\!\mathrm{cm}\) سيكون متوافقاً مع الفرضية؛ حد أقل من \(10^{-28}\;e\!\cdot\!\mathrm{cm}\) سيدحضها.
		\item \textbf{بحوث الأكسيون (ADMX, MADMAX, IAXO).}
		إذا تحققت آلية بيكي-كوين، فيجب أن تكون كتلة الأكسيون التي تنبأت بها YUCT قابلة للحساب بمجرد تقييم \(f_a\) من \(L_\theta\). بحث مخصص في نافذة الكتلة المقابلة سيختبر الفرضية.
		\item \textbf{آثار كونية.}
		إذا كان تداخل النظام-الفوضى أساسياً، فقد يؤثر أيضاً على تحولات الطور في الكون المبكر، تاركاً بصمات محددة في الخلفية الكونية الميكروية أو في توزيع المادة المظلمة. يمكن تقييد هذه باستخدام Planck و Euclid وبيانات CMB‑S4 المستقبلية.
	\end{enumerate}
	
	\section{دعوة للتعاون}
	\label{sec:invite}
	
	إطار YUCT، بما في ذلك لاغرانجيه الكامل، وجميع الملاحق، وقاعدة البيانات الواسعة لأعماق التنسيق، متاح مفتوحاً على Zenodo~\cite{Y1,Y2}. ندعو الفيزيائيين والرياضيين وعلماء الحاسوب للانضمام إلى الجهد لاشتقاق أو اختبار أو دحض الفرضية المقدمة هنا. المؤلف (أليكسي ياكوشيف، \href{mailto:alexey@yakushev.eu}{alexey@yakushev.eu}) يرحب بالمراسلات والاقتراحات ومقترحات التعاون.
	
	\section{الاستنتاج}
	\label{sec:conclusion}
	
	لقد اقترحنا صيغة جبرية محددة خالية من البارامترات للعمق الفعال لمعامل QCD \(\theta\) ضمن نظرية التنسيق الموحدة ياكوشيف:
	\[
	L_\theta = \frac{1}{2}L_W + \frac{S_{\mathrm{odd}}}{S_{\mathrm{even}}} = 49.5,
	\]
	مما يعطي \(\theta \approx 2.3 \times 10^{-10}\). هذه الفرضية بسيطة رياضياً، ومدفوعة فيزيائياً، وقابلة للدحض المباشر من خلال تجارب nEDM والأكسيون القادمة. تحل محل المحاولات السابقة التي اعتمدت على تعيينات أعداد صحيحة عشوائية وتفتح طريقاً واضحاً نحو حل كامل لمشكلة CP القوية ضمن YUCT. ما إذا كانت الفرضية ستصمد أمام التدقيق التجريبي يبقى أن نرى؛ قيمتها تكمن في تقديم تخمين دقيق قابل للاختبار يمكنه توجيه البحث المستقبلي.
	
	\vspace{1cm}
	\noindent\textbf{التاريخ:} يونيو 2026\\
	\textbf{الإصدار:} 2.0\\
	\textbf{الحالة:} فرضية مع تنبؤ جبري محدد؛ مفتوحة للاختبار التجريبي والاشتقاق النظري.
	
	\newpage
	\begin{thebibliography}{99}
		\bibitem{Y1} A.V. Yakushev, \textit{Yakushev Unified Coordination Theory (YUCT)}, Zenodo, 2026. DOI:10.5281/zenodo.18444599.
		\bibitem{Y2} A.V. Yakushev, Appendix AF: The Algebraic Framework of Reality in YUCT, 2026.
		\bibitem{Y3} A.V. Yakushev, Appendix \(\Lambda\): Resolution of the Hierarchy and Cosmological Constant Problems, 2026.
		\bibitem{Y4} A.V. Yakushev, Appendix J: Complete Derivation of Standard Model Flavor Parameters, 2026.
		\bibitem{Feig} M.J. Feigenbaum, \textit{J. Stat. Phys.} \textbf{19}, 25 (1978).
		\bibitem{nEDM2020} C. Abel et al. (nEDM Collaboration), \textit{Phys. Rev. Lett.} \textbf{124}, 081803 (2020).
	\end{thebibliography}
	
\end{document}